\documentclass[../main.tex]{subfiles}
\graphicspath{{\subfix{../images/}}}
\begin{document}

% ------------------------------------------------------------
% Selbstständig: funktioniert auch ohne Zusatz im Preamble der main.tex
% ------------------------------------------------------------
\usetikzlibrary{shapes.multipart,shapes.geometric}   % für rectangle split / ellipse

% Querformat-Seite nur definieren, falls main.tex sie nicht schon kennt
\ifcsname querseite\endcsname\else
  \newenvironment{querseite}{%
    \clearpage
    \newgeometry{margin=1cm,top=1.5cm,bottom=1.5cm}%
    \begin{landscape}%
    \thispagestyle{empty}%
  }{%
    \end{landscape}%
    \clearpage
    \restoregeometry
  }
\fi

%\begin{querseite}

% Kapitelüberschrift steht IN der Querformat-Seite, damit sie mit der
% Mindmap auf derselben Seite landet. In der main.tex darf dafür KEIN
% eigenes \section{Mindmap} mehr vor dem \subfile stehen.
\section{Mindmap}
\label{Brainstorming}

% TODO: Kurz erklären, wie die Ideen entstanden sind.

\definecolor{bshaupt}{RGB}{40,90,188}
\definecolor{bshell}{RGB}{232,241,250}

\tikzset{
  bszentrum/.style={ellipse, draw=bshaupt, fill=bshaupt, text=white, very thick,
    minimum width=5.4cm, minimum height=3cm, text width=4.2cm,
    align=center, font=\bfseries\Large},
  bsast/.style={rectangle split, rectangle split parts=2, draw=bshaupt, thick,
    rounded corners=3pt, text width=6.6cm, inner sep=5pt,
    rectangle split part fill={bshaupt,bshell}, align=left, font=\small},
  bsverbindung/.style={draw=bshaupt, very thick}
}

\providecommand{\bszweig}[4]{% Name, Position, Titel, Unterpunkte
  \node[bsast] (#1) at (#2) {%
    \textcolor{white}{\bfseries #3}
    \nodepart{second}
    \begin{minipage}{6.4cm}
      \begin{itemize}[leftmargin=1.1em,itemsep=1pt,topsep=1pt,parsep=0pt,label=\textcolor{bshaupt}{$\bullet$}]
        #4
      \end{itemize}
    \end{minipage}};
}

\vspace*{\fill}  % Grafik vertikal zentrieren

\centering
\ifcsname singlespacing\endcsname\singlespacing\fi  % Knoten nicht durch 1,5-fachen Zeilenabstand aufblähen
\adjustbox{max width=\linewidth,max totalheight=0.78\textheight,center}{%
\begin{tikzpicture}[x=1cm,y=1cm]

\node[bszentrum] (Z) at (0,0) {Brainstorming\\Gehäusekonzept};

% Linke Seite
\bszweig{A}{-8.6,7.0}{1\; Befestigung der Platine im Gehäuse}{
  \item Gar nicht
  \item Schrauben ins Material
  \item Einschmelzgewinde
  \item Kleben
  \item Verguss (Epoxidharz)}

\bszweig{B}{-8.6,2.5}{2\; Befestigung der Komponenten}{
  \item Steckverbindungen (Pins)
  \item Löten (SMD / Through-Hole)
  \item Schraubverbindung
  \item Kleben}

\bszweig{C}{-8.6,-2.5}{3\; Energieversorgung / Wandlung (24\,V / 5\,V)}{
  \item Einzelne Netzteile für 24\,V und 5\,V
  \item Buck-/Boost-Converter (24\,V $\rightarrow$ 5\,V)
  \item Trafo mit Gleichrichter (230\,V-Anschluss)
  \item Zener-Diode / Spannungsteiler für 5\,V
  \item Akku oder Netzversorgung}

\bszweig{D}{-8.6,-7.0}{4\; Befestigung der Motorkabel}{
  \item Steckverbindung
  \item Anlöten
  \item Schraubklemmen
  \item Magnetpins
  \item Keine (fest verdrahtet)}

% Rechte Seite
\bszweig{E}{8.6,7.0}{5\; Wärmeabfuhr}{
  \item Keine
  \item Natürliche Belüftung (Lüftungsschlitze)
  \item Kühlrippen
  \item Lüfter
  \item Wasserkühlung
  \item Flüssiger Stickstoff (LN$_2$)}

\bszweig{F}{8.6,2.5}{6\; Fertigung des Gehäuses}{
  \item 3D-Druck
  \item CNC
  \item Manuelles Spanen
  \item Spritzgussverfahren}

\bszweig{G}{8.6,-2.5}{7\; Verbindung Gehäuse und Roboter}{
  \item Steckverbindung / Klemmen
  \item Magnetisch
  \item Schraubverbindung
  \item Kleben
  \item Einteilig}

\bszweig{H}{8.6,-7.0}{8\; Material des Gehäuses}{
  \item PLA
  \item PETG
  \item ABS
  \item Aluminium
  \item Unlegierter Baustahl (S235JR)
  \item Acryl (durchsichtig)}

% Geschwungene Verbindungen, verteilt über die Ellipse
\draw[bsverbindung] (Z.150) to[out=180,in=0] (A.east);
\draw[bsverbindung] (Z.172) to[out=180,in=0] (B.east);
\draw[bsverbindung] (Z.188) to[out=180,in=0] (C.east);
\draw[bsverbindung] (Z.210) to[out=180,in=0] (D.east);

\draw[bsverbindung] (Z.30)  to[out=0,in=180] (E.west);
\draw[bsverbindung] (Z.8)   to[out=0,in=180] (F.west);
\draw[bsverbindung] (Z.-8)  to[out=0,in=180] (G.west);
\draw[bsverbindung] (Z.-30) to[out=0,in=180] (H.west);

\end{tikzpicture}}

\captionof{figure}{Brainstorming: Ideensammlung zum Gehäusekonzept}
\label{fig:brainstorming}

\vspace*{\fill}
%\end{querseite}

\end{document}